% !TeX root = ../main.tex
% !TeX encoding = UTF-8

\part{The Protocol}

\chapter{The Digital Alibi Assessment, What Forensic Reconstructibility Actually Requires}

The examiner places the Decision Reconstruction Report on the table. Twenty-five material AI-assisted decisions, five credit decisions, five underwriting decisions, five hiring decisions, five procurement awards, five regulatory compliance determinations, examined against the four Digital Alibi elements. For twenty of the twenty-five, the Information Picture is a total black hole. No contemporaneous record of what the AI system knew. No record of what it was not given. For fourteen of the twenty-five, the Accountability Chain breaks before it reaches the board. For eleven, the Human Oversight Record shows a sign-off that could not have represented genuine review at the volume and speed at which it was processed. The board chair asks the question that will be asked of every board in every regulated organisation within the next eighteen months: how did we not know this?

The answer is that the assessment is the first time anyone looked. The board had approved the systems. They had approved the autonomy boundaries, the permission sets, under which those systems operated. But they had never confirmed that the standing authorisation they granted was being executed with any evidence of position. Every action the systems took was pre-authorised. None of it was recorded at the output-moment level in a form that survives scrutiny. The Digital Alibi gap exists in every organisation that has not confirmed its governance address with forensic precision. The assessment does not create the gaps it finds. It names existing exposures in a form that allows a deliberate governance response rather than an unexpected regulatory finding. That is the difference between conducting the assessment now and waiting for a regulator to conduct it for you.

Part Four is where the board acts. Not plans to act. Not commissions a review. Acts: specifically, in sequence, with named individuals accountable for each stage, against a governance architecture that can be confirmed by an independent examiner. The Digital Alibi Assessment is the first act.

\scenebreak

\section{What the Assessment Is, and What It Is Not}

\section{\texorpdfstring{\emph{It tells the board exactly what evidence existed at the moment a decision was made, or precisely why that evidence is missing.}}{It tells the board exactly what evidence existed at the moment a decision was made, or precisely why that evidence is missing.}}

The Digital Alibi Assessment is a structured protocol for establishing where a board currently stands against the four Digital Alibi elements across its most recent material AI-assisted decisions. It produces a Decision Reconstruction Report, a forensic account of what existed, what was retrievable, and what was missing at the moment of each decision examined. The confirmed governance address that emerges from the assessment is the starting point for the Protocol that the remaining chapters of Part Four describe.

The assessment is not a technical audit. Technical audits examine systems for compliance with technical standards. The Digital Alibi Assessment examines the evidentiary record that exists for specific decisions, whether the information picture at each decision moment was captured, whether it is retrievable, and whether the accountability chain that runs from system to board is documented with decision-moment precision. These are governance questions, not technical ones. The answers may require technical investigation to establish, but the questions are asked on behalf of the board, not on behalf of the technology function.

The assessment is not a self-assessment. Chapter 14 established that the gap between self-assessed and confirmed governance address is itself a governance finding. Self-assessment produces the address the organisation believes it has. The Digital Alibi Assessment produces the address it can demonstrate. An organisation that self-assesses at A3·E3·C3·R3·M3 and is confirmed at A2·E2·C2·R2·M2 has learned, as the assessment's first deliverable, that its governance confidence was not founded in evidence. That recalibration is worth the assessment before any other finding is considered.

\begin{quote}
\emph{The Digital Alibi Assessment is not a technical audit and not a self-assessment. It is a forensic examination of the evidentiary record that exists for specific decisions. The questions are governance questions. The examiner is independent. The findings are board-level governance decisions, not technical remediation tasks.}
\end{quote}

\scenebreak

\section{Who Conducts the Assessment and Who Owns the Findings}

The assessment must be conducted independently of the teams whose governance it examines. This is the independence design principle from Chapter 5 applied at the governance level: the audit of the Digital Alibi architecture must be independent of the parties whose accountability it documents. An assessment conducted by the technology team on their own model versioning records, or by the risk function on their own oversight records, does not produce an independent finding. It produces a self-attestation.

\begin{quote}
\emph{Independence is not a procedural nicety. It is the only condition under which a board can trust the forensic findings. An examiner with institutional interest in a favourable outcome has already shaped the finding before the report is written.}
\end{quote}

Independence has two components. The examiner must be independent of the systems and processes being examined. And the examiner must be independent of the institutional interest in a favourable outcome. An internal audit function that is structurally independent of the technology team meets the first component. Whether it meets the second depends on the institution's governance culture, whether internal audit is genuinely empowered to produce findings that are uncomfortable for the board, or whether institutional dynamics shape the findings before they are presented.

The assessment should be owned by a named non-executive director. Not the Chief Risk Officer, who may have institutional interest in the findings. Not the CTO, who is accountable for the systems being examined. A non-executive director whose accountability is to the board as a whole, and who has no departmental interest in a particular outcome, is the appropriate owner. The NED who owns the assessment presents its findings to the audit committee or board risk committee, and proposes the governance response that the findings require.

\begin{longtable}[]{@{}
  >{\raggedright\arraybackslash}p{(\columnwidth - 0\tabcolsep) * \real{1.0000}}@{}}
\toprule\noalign{}
\begin{minipage}[b]{\linewidth}\raggedright
\textbf{BOARD NOTE, NED OWNERSHIP OF THE DIGITAL ALIBI ASSESSMENT}
\end{minipage} \\
\midrule\noalign{}
\endhead
\bottomrule\noalign{}
\endlastfoot
The NED appointed to own the assessment must have no operational stake in the AI systems under review. This is not a formality, it is the condition on which the assessment remains a governance check rather than a departmental defence. \\
The NED owner should: confirm the scope of the assessment before it begins; receive the Decision Reconstruction Report directly from the examiner, not filtered through the technology or risk function; present the findings to the audit committee with their own analysis of the governance response required; and accept named accountability for the board's response to each gap identified. \\
A NED who delegates ownership of the assessment to an executive function has delegated the governance check to the function it is checking. The independence principle does not survive that delegation. \\
\end{longtable}

The size and formality of the assessment is not fixed. For a board that is conducting its first Digital Alibi Assessment, the protocol described in this chapter represents a full implementation. For a board that has conducted previous assessments and is conducting a periodic review, a targeted examination of specific decision types or specific systems may be appropriate. The principle that governs the scope is the materiality test from Chapter 12: the assessment should examine the decisions where the Digital Alibi gap creates the most significant regulatory, legal, or human-consequences risk.

\scenebreak

\section{The Five Decision Types}

The Digital Alibi Assessment examines five material decision types that represent the highest concentration of Digital Alibi risk in most regulated organisations. For each decision type, the five most recent material decisions are examined against the four Digital Alibi elements. Twenty-five decision examinations in total, each producing a partial or complete Decision Reconstruction Record that becomes the foundation of the Decision Reconstruction Report.

The five decision types are:

Credit and lending decisions. AI-assisted decisions that determine whether an individual or organisation receives credit, at what price, and on what terms. This includes consumer credit scoring, mortgage affordability assessment, business lending decisions, and any AI-assisted pricing or terms-setting that affects borrowers. These decisions engage FCA Consumer Duty obligations, GDPR Article 22 automated decision rights, and the EU AI Act high-risk classification.

Insurance underwriting and claims. AI-assisted decisions that determine whether an individual or organisation is insurable, at what premium, and on what coverage terms, and AI-assisted decisions that determine whether a claim is valid, at what settlement value, and whether to pay or decline. These decisions engage FCA and PRA conduct obligations, Lloyd's market governance requirements for syndicates, and EU AI Act high-risk classification for certain underwriting categories.

Hiring and people decisions. AI-assisted decisions that affect whether an individual is selected for employment, promotion, or termination, including CV screening, candidate ranking, interview scheduling, and performance assessment systems. These decisions engage Equality Act 2010 obligations in the UK, EU AI Act high-risk classification for employment AI, and ICO guidance on AI in HR.

Procurement and contract award. AI-assisted decisions that determine which supplier is selected, which bid is accepted, and which contract is awarded. These decisions engage public procurement law for public sector organisations, Equality Act obligations for private sector organisations with supplier diversity commitments, and competition law obligations in certain sectors.

Regulatory compliance determinations. AI-assisted decisions that determine whether a regulatory obligation has been met, including transaction monitoring for AML, sanctions screening, conduct surveillance, and market abuse detection. These decisions engage the most direct personal liability for named Senior Managers under SM\&CR, because regulatory compliance failures in these categories are the ones most frequently subject to FCA enforcement action.

\begin{longtable}[]{@{}
  >{\raggedright\arraybackslash}p{(\columnwidth - 0\tabcolsep) * \real{1.0000}}@{}}
\toprule\noalign{}
\begin{minipage}[b]{\linewidth}\raggedright
\textbf{FOR EACH OF THE FIVE DECISION TYPES, FIVE QUESTIONS}
\end{minipage} \\
\midrule\noalign{}
\endhead
\bottomrule\noalign{}
\endlastfoot
1. Can we identify the specific AI system, at the specific version deployed, that contributed to each of the five most recent material decisions in this category? \\
2. Can we produce the complete information picture, what was there, what was retrievable, what was missing, for each of those five decisions, from the contemporaneous record rather than from current system states? \\
3. Is there a contemporaneous human oversight record for each decision, naming the reviewer, what they were shown, what their override authority was, and what they decided? \\
4. Is the accountability chain for each decision documented at the output-moment level, naming the model owner, the decision process owner, and the board-level risk acceptance authority who were current at the specific decision moment? \\
5. Can all four of the above be produced within five working days, from the contemporaneous record, without manual reconstruction? \\
\end{longtable}

A board that receives affirmative answers to all five questions for all five decision types has a Digital Alibi architecture that is functioning at the output-moment level. A board that receives a negative answer to any question for any decision type has identified a specific, named Digital Alibi gap at a specific decision type. The gap is the governance finding. The Protocol is the response.

\scenebreak

\section{The Assessment in Four Phases}

\section{Phase 1 AI Systems Inventory Confirmation}

Before any decision can be examined, the board must confirm which AI systems contributed to it. Phase 1 identifies every AI system contributing to decisions in the five material decision types, including vendor-embedded AI and shadow AI absent from the formal inventory.

The inventory confirmation begins with the formal AI inventory and expands it systematically. For each decision examined, the core questions are: which AI system produced the output that informed this decision, at which version, and is that version separately retrievable if the model has since been updated? Are there upstream AI systems, data enrichment, feature engineering pipelines, data quality assessment tools, that contributed to the inputs the primary model received? These upstream systems are material to the Information Picture even when they are not in the formal AI inventory.

Phase 1 produces a confirmed AI systems inventory for each decision type: every AI system, at the version deployed at each decision, that contributed to the decisions being examined. This inventory is the Decision Identity element of the Digital Alibi for each decision. If it cannot be produced for any decision, the assessment has found its first gap before examining any other element.

\section{Phase 2 Decision Reconstruction}

Phase 2 is the forensic examination of each decision against the four Digital Alibi elements. For each of the twenty-five decisions examined, five decisions across each of the five decision types, the assessment produces a Decision Reconstruction Record: a structured account of what existed, what was retrievable, and what was missing at the decision moment for each of the four elements.

The Decision Reconstruction Record is not a report of what the organisation intended to capture or what its governance framework specifies should be captured. It is an account of what actually exists in the contemporaneous record for this specific decision. The examiner attempts to retrieve the Decision Identity, The Information Picture, The Human Oversight Record, and The Accountability Chain from the contemporaneous record, from the logs, the version control systems, the data lineage records, the oversight audit trails, and the accountability documentation as they existed at the decision moment.

What cannot be retrieved from the contemporaneous record is a gap. Not a gap in the governance framework, the framework may specify that it should have been captured. A gap in the governance record: the evidence that the governance operated at this decision, at this moment, with this information picture, under this accountability chain. That gap is the Digital Alibi gap, named and specific, for this decision.

\section{Phase 3 Governance Address Confirmation}

Phase 3 uses the findings from the twenty-five Decision Reconstruction Records to confirm the organisation's governance address across A·E·C·R·M with forensic precision. The confirmed address reflects what the evidence shows, not what the organisation believes or aspires to.

The gap between the self-assessed address from Chapter 14's diagnostic and the confirmed address from the assessment findings is itself a governance deliverable. An organisation that self-assessed at A3·E3·C3·R3·M3 and finds, from the Decision Reconstruction Records, that accountability chains are not documented at the decision-moment level for any of the twenty-five decisions examined, that shadow AI contributed to decisions in three of the five decision types without appearing in the formal inventory, and that the human oversight records show performative rather than genuine review capacity, has confirmed its address at A2·E2·C2·R2·M2 regardless of its governance framework documentation.

The confirmed address maps directly to the sector ceiling established in Chapter 14. The gap between the two, on each facet independently, is the governance investment required to reach defensible compliance.

\section{Phase 4 Board Presentation and Protocol Initiation}

Phase 4 is the presentation of the Decision Reconstruction Report to the audit committee or board risk committee, and the board's formal adoption of a Protocol response. The findings are board-level governance decisions requiring named accountability, not IT remediation tasks.

The board's response to each gap takes one of three forms. For accountability chain gaps: a named individual accepts accountability for the gap and commits to a specific date by which the decision-moment attribution architecture is operational. For information picture gaps: the board accepts the liability position for historical decisions and commits to the technical architecture that will capture the information picture going forward. For human oversight record gaps: the board restructures the oversight process to create genuine review capacity and names the individual accountable for confirming that capacity is genuine, not performative.

Each response requires a named individual and a named date. A response that does not name both has not initiated the Protocol. It has noted the finding.

\begin{longtable}[]{@{}
  >{\raggedright\arraybackslash}p{(\columnwidth - 0\tabcolsep) * \real{1.0000}}@{}}
\toprule\noalign{}
\begin{minipage}[b]{\linewidth}\raggedright
\textbf{THE BOARD CHECKLIST, DIGITAL ALIBI ASSESSMENT READINESS}
\end{minipage} \\
\midrule\noalign{}
\endhead
\bottomrule\noalign{}
\endlastfoot
1. Has the organisation produced a current inventory of all AI systems contributing to material decisions, confirmed to include vendor-embedded and shadow AI? \\
2. Is there a named non-executive director who owns the Digital Alibi Assessment and will present its findings to the audit committee? \\
3. Are model versioning, parameter logging, and data lineage capture operational for all material AI systems, at output-moment level, not system level? \\
4. Does the human oversight record for each material AI system show genuine review capacity, or does the volume, speed, and format of review make genuine oversight structurally impossible? \\
5. Has the board received a formal Decision Reconstruction Report in the last twelve months, not a governance framework review, but a forensic account of what existed at the moment of specific past decisions? \\
\end{longtable}

\scenebreak

\section{The Assessment Is the Beginning, Not the End}

The Digital Alibi Assessment produces a finding. It does not close the gap. The confirmed address it produces, the forensically precise location of the organisation in the governance landscape, is the starting point for the Protocol chapters that follow. The Decision Reconstruction Report identifies where the evidence was missing, which decisions are most exposed, and which facets of the governance address require the most urgent investment.

The assessment does not produce liability. The gaps it identifies are not new exposures, they are existing exposures that have been named and quantified. A board that conducts the assessment and finds significant Digital Alibi gaps has not created those gaps by examining them. It has discovered them in a form that allows a deliberate governance response rather than an unexpected regulatory finding.

The most common question boards ask after receiving their first Decision Reconstruction Report is what to do first. The answer from the confirmed governance address is clear: address the facet on which the gap between the confirmed address and the sector ceiling is widest, because that gap is the one on which the regulatory and legal exposure is most concentrated. For most regulated organisations, that facet is A, Accountability. The accountability chain that runs from the AI system to the named board-level individual is the element most consistently absent from the decision-moment record, and the element on which SM\&CR personal liability most directly depends.

\begin{quote}
\emph{The assessment does not create the gaps it finds. It names existing exposures in a form that allows a deliberate governance response. A board that discovers significant gaps has not failed the Digital Alibi standard at the moment of discovery. It has begun to meet it.}
\end{quote}

Chapter 17 describes the Decision Map, the operational architecture that turns the assessment's findings into a prospective governance record that captures the complete information picture at every future material decision. Building the Decision Map is the first act of Protocol implementation. It transforms the assessment's historical finding into a forward-looking governance infrastructure.
